%% ============================================================
%%  BOOK 2 — CHAPTER 9: MARTIAN COLONY ARCHITECTURE
%%  AND THE DM3 HABITABILITY CRITERION
%%  Pablo Nogueira Grossi · G6 LLC · Newark NJ · 2026
%% ============================================================

\chapter{Martian Colony Architecture and the dm3 Habitability Criterion}
\label{ch:mars}

\begin{center}
\itshape
``A planetary environment can support life\\
if and only if the ambient noise level $\sigma$ satisfies\\
$\sigma < \tau$ for at least one class of oscillatory\\
system present in that environment.''\\[0.5em]
\normalfont\small --- The dm3 Habitability Criterion, derived from Theorem~A
\end{center}

\vspace{0.5cm}

\noindent The question of whether Mars can support life is one of the
oldest questions in planetary science and one of the least precisely
stated. What does ``support life'' mean, mathematically? What is the
threshold? What is being measured? What constitutes a falsifiable test?

The dm3 framework answers these questions for the first time with a
quantitative, substrate-agnostic, mathematically derived criterion.
The criterion does not depend on what life is made of. It depends on
whether the ambient noise environment of a planet permits the existence
of a stable oscillatory attractor --- a limit cycle --- in at least
one physical system present on that planet.

This chapter derives the criterion, applies it to Mars, derives three
falsifiable predictions testable with existing or near-term instrumentation,
and then extends the framework to the architectural geometry of a
permanent Martian colony.

%% ─────────────────────────────────────────────────────────────
\section{The dm3 Habitability Criterion}
\label{sec:mars-criterion}
%% ─────────────────────────────────────────────────────────────

The central result of Volume~II (Generative Contact Mechanics, Theorem~A)
is the following: a dm3 system $\mathfrak{D} = (S, g, Y, \Gamma, V, \sigma)$
admits a unique exponentially stable limit cycle $\Gamma$ with stochastic
stability below the embodiment threshold $\tau = \sqrt{c/\kappa}$, where
$c > 0$ is the Lyapunov contraction rate and $\kappa$ is the noise
coupling constant.

Below $\tau$: noise reinforces the orbit. The system is alive in the
dm3 sense --- it maintains stable oscillation against perturbation.
Above $\tau$: noise disrupts the orbit. The attractor collapses.

\begin{definition}[dm3 Habitability Criterion]
\label{def:habitability}
A planetary environment $\mathcal{E}$ is \emph{dm3-habitable} if there
exists at least one physical system $\mathfrak{D}$ present in $\mathcal{E}$
such that
\[
\sigma_{\mathcal{E}}(\mathfrak{D}) < \tau(\mathfrak{D}),
\]
where $\sigma_{\mathcal{E}}(\mathfrak{D})$ is the ambient noise amplitude
experienced by $\mathfrak{D}$ in $\mathcal{E}$, and
$\tau(\mathfrak{D}) = \sqrt{c(\mathfrak{D})/\kappa(\mathfrak{D})}$
is the embodiment threshold of $\mathfrak{D}$.
\end{definition}

This is the first mathematically precise, substrate-agnostic, quantitative
threshold for the possibility of life on any world. It does not ask
whether the right molecules are present. It asks whether the physical
environment permits stable oscillatory structure to exist at all. Life,
in the dm3 sense, is what happens when a system crosses below its
embodiment threshold and locks onto a limit cycle.

Applied to Earth: every biological oscillator we know --- HPA stress
cycles, neural rhythms, circadian clocks, immune adaptation cycles ---
satisfies $\sigma < \tau$ with explicit, measurable values. The
canonical dm3 triple $(T^*, \mu_{\max}, \tau) = (2\pi, -2, 2)$
characterizes the class of systems that biological evolution has
selected for. Earth is dm3-habitable. This is not a tautology. It is
a verification of the criterion against known biology.

Applied to Mars: the question is whether Mars has any physical system
for which $\sigma_{\text{Mars}} < \tau$. We identify four candidate
systems and derive three falsifiable predictions.

%% ─────────────────────────────────────────────────────────────
\section{Mars Physical Parameters}
\label{sec:mars-params}
%% ─────────────────────────────────────────────────────────────

The relevant physical parameters of the Martian environment are:

\begin{table}[h]
\centering
\small
\begin{tabular}{lll}
\toprule
\textbf{Parameter} & \textbf{Value} & \textbf{Source} \\
\midrule
Surface gravity & $g_\Mars = 3.72\,\mathrm{m/s^2}$
  & $= 0.379\,g_\oplus$ \\
Surface pressure & $P = 610\,\mathrm{Pa}$
  & $= 0.60\%$ of Earth \\
Mean surface temperature & $T = 210\,\mathrm{K}$ ($-63^\circ$C) & \\
Sol period & $T_{\text{sol}} = 88{,}642\,\mathrm{s}$ & $= 24.62$ hours \\
Speed of sound (CO$_2$) & $c_s = 235.7\,\mathrm{m/s}$ & \\
Thermal diffusivity (regolith) & $\kappa_T \approx 10^{-7}\,\mathrm{m^2/s}$ & \\
Thermal skin depth & $\delta = \sqrt{2\kappa_T/\omega_{\text{sol}}}
  = 5.31\,\mathrm{cm}$ & \\
\bottomrule
\end{tabular}
\caption{Mars physical parameters relevant to the dm3 habitability analysis.}
\label{tab:mars-params}
\end{table}

The skin depth of $5.31\,\mathrm{cm}$ is the depth at which the
daily thermal wave decays to $1/e$ of its surface amplitude. Below
approximately $3$--$5$ skin depths ($\approx 15$--$27\,\mathrm{cm}$),
the thermal environment is effectively decoupled from the violent
diurnal surface cycle and becomes a stable oscillatory medium.

%% ─────────────────────────────────────────────────────────────
\section{Four Candidate dm3 Systems on Mars}
\label{sec:mars-candidates}
%% ─────────────────────────────────────────────────────────────

\subsection{System 1: Subsurface Thermal Oscillator}

The Martian sol imposes a thermal forcing on the surface and subsurface
at period $T_{\text{sol}} = 88{,}642\,\mathrm{s}$. Below the skin depth,
this forcing drives a damped oscillation in temperature --- a physical
system with a well-defined period, exponential damping (controlled by
thermal diffusivity), and stochastic perturbation from dust storm
insolation variability.

\textbf{dm3 identification:}
\begin{itemize}
\item $T^* = T_{\text{sol}} = 88{,}642\,\mathrm{s}$ (period of limit cycle)
\item $\mu_{\max} = -1/\delta^2 \cdot \kappa_T < 0$ (exponential stability,
  rate set by thermal diffusivity)
\item $\sigma = $ insolation variability from dust storms
\item $\tau = \sqrt{c_\text{thermal}/\kappa_\text{thermal}}$
\end{itemize}

The key question: does Mars dust storm variability push $\sigma$ above
$\tau$ at depths accessible to potential biology ($> 5\,\mathrm{cm}$)?
The InSight mission's Heat Flow and Physical Properties Package (HP3)
was designed to measure exactly this. The data, analyzed through the
dm3 framework, provides a direct test of the habitability criterion
for the thermal oscillator.

\subsection{System 2: Dust Devil Pressure Oscillations}

Martian dust devils are convective vortices with core pressure drops
of $0.5$--$4\,\mathrm{Pa}$ and vortex frequencies of approximately
$0.1\,\mathrm{Hz}$. They are not random noise. They are structured
pressure oscillations with a well-defined frequency, amplitude, and
decay time. They are candidate dm3 systems.

\textbf{dm3 identification:}
\begin{itemize}
\item $T^* \approx 10\,\mathrm{s}$ (vortex period)
\item $\mu_{\max}$: set by vortex decay time (seconds to minutes)
\item $\sigma$: ambient pressure noise from the background atmosphere
\item $\tau$: threshold below which the vortex maintains coherence
\end{itemize}

The noise fraction of surface pressure is $\sigma/P \approx 0.008$
for typical ambient fluctuations. The dust devil maintains its
structure against this background noise. Its Arnold tongue coupling
to the planetary atmospheric modes (via $\mathcal{A}_{1:1}$ at the
sol frequency) is the testable signature: a pressure sensor array
should detect mode-locking between dust devil vortex frequency and
the $1/T_{\text{sol}}$ planetary mode.

\subsection{System 3: Perchlorate Brine Chemistry}

Perchlorate brines on Mars remain liquid at temperatures above
approximately $200\,\mathrm{K}$ --- conditions that exist seasonally
at mid-latitudes and possibly year-round in the deep subsurface.
Reaction-diffusion chemical systems in brines can sustain limit-cycle
oscillations (Belousov-Zhabotinsky type). If Martian perchlorate brine
chemistry admits such oscillations, the dm3 habitability criterion
becomes:
\[
\sigma_{\text{electrochemical}} < \tau_{\text{brine chemistry}}.
\]
The noise amplitude is set by thermal fluctuations in the brine;
the threshold is set by the reaction rates and diffusivity. This
is measurable with in-situ electrochemical sensing at brine-bearing sites.

\subsection{System 4: Atmospheric Cavity Resonance}

The Martian CO$_2$ atmosphere has a planetary acoustic resonance:
\[
f_{\text{Helmholtz}} = \frac{c_s}{2\pi R_\Mars}
= \frac{235.7}{2\pi \times 3.39 \times 10^6}
= 0.011\,\mathrm{mHz},
\]
with period $T \approx 25.1$ hours --- close to the sol period.
The near-coincidence of the acoustic cavity period and the sol period
creates conditions for Arnold tongue coupling: the thermal forcing at
$T_{\text{sol}} = 24.62\,\mathrm{hours}$ drives the atmospheric
cavity at a frequency close to its natural resonance. This is a
planetary-scale $k$:$m$ resonance, with the dm3 stability radius
$\varepsilon_0 = 1/3$ governing whether the coupling is structurally
stable against seasonal variability.

%% ─────────────────────────────────────────────────────────────
\section{Three Falsifiable Predictions}
\label{sec:mars-predictions}
%% ─────────────────────────────────────────────────────────────

\begin{enumerate}

\item \textbf{Thermal habitability depth.}
The subsurface thermal oscillator at depth $d > 3\delta \approx 16\,\mathrm{cm}$
maintains $\sigma < \tau$ throughout the Martian year, including global
dust storm season. \emph{Test:} reanalysis of InSight HP3 thermal data
at depth; new penetrometer missions targeting $20$--$50\,\mathrm{cm}$
depth with continuous temperature logging. The dm3 criterion predicts
a sharp transition in oscillation coherence at the thermal skin depth
boundary.

\item \textbf{Brine electrochemical oscillation.}
Perchlorate brine at temperature $T > 200\,\mathrm{K}$ exhibits
chemical limit-cycle oscillations with embodiment threshold $\tau_{\text{brine}}$
measurable via electrochemical impedance spectroscopy. \emph{Test:}
in-situ electrochemical sensor at brine-candidate sites (Recurring
Slope Lineae, subsurface ice-brine interfaces). The prediction is
that the noise floor of Martian electrochemical environments is below
$\tau_{\text{brine}}$ at these sites.

\item \textbf{Dust devil Arnold tongue signature.}
Martian dust devil vortices exhibit mode-locking to the planetary
atmospheric background via Arnold tongue $\mathcal{A}_{1:1}$, with
the coupling structurally stable within $\varepsilon_0 = 1/3$.
\emph{Test:} barometric pressure array (minimum 3 sensors, spacing
$\sim 1\,\mathrm{km}$) deployed at high-dust-devil-frequency sites
(e.g., Gusev Crater, Jezero Crater rim). The prediction is that
dust devil occurrence is correlated with sol phase at the
$\mathcal{A}_{1:1}$ locking fraction, not uniformly distributed.

\end{enumerate}

These predictions are derived from the dm3 framework, not from
biological assumptions. They do not require the presence of organisms.
They test whether the physical environment satisfies the mathematical
conditions for stable oscillation --- the necessary precondition
for life in the dm3 sense.

%% ─────────────────────────────────────────────────────────────
\section{Closed-Loop Life Support as dm3 Operator Chain}
\label{sec:mars-lifesupport}
%% ─────────────────────────────────────────────────────────────

A Martian colony is itself a dm3 system. Its life support architecture
must maintain stable oscillatory attractors --- resource cycles,
atmospheric cycles, biological cycles --- against the noise of the
Martian environment. The operator chain $G = U \circ F \circ K \circ C$
governs how those cycles are established and maintained.

\textbf{C (Compression):} Resource compression. The colony reduces
all life support to a minimal set of coupled cycles: $\mathrm{CO_2}
\to \mathrm{O_2}$ conversion, water recycling, food production, waste
processing. Degrees of freedom are compressed to the minimum viable
set. This is the first operator application.

\textbf{K (Curvature):} Stress accumulation. As the colony grows,
physiological and engineering loads accumulate. Crew stress,
equipment wear, resource depletion, radiation exposure --- these drive
the system toward critical thresholds. The curvature operator $K$
models this accumulation: the system is driven toward $\kappa^*$,
the fold threshold.

\textbf{F (Fold):} System fold events. Atmosphere breach, crop failure,
equipment cascade failure, crew medical emergency --- these are
fold events in the dm3 sense: rank-1 Jacobian losses at which the
system either transitions to a new branch or collapses to $B_3$
(the collapse boundary). The fold is not the end. It is the transition.

\textbf{U (Unfold):} Recovery branch selection. Redundant system
activation, biological adaptation, ISRU rerouting, crew resilience ---
these are the unfold operator: gradient descent onto a new stable
branch. The allostatic unification law (Theorem~B) governs coupled
subsystems:
\[
\tau_{12} \leq \min(\tau_1, \tau_2).
\]
When the atmospheric system (system 1) and the crew physiology
(system 2) are coupled under stress, the combined system's noise
tolerance is bounded above by the weaker of the two. The colony
architect must design so that $\tau_{\text{atmo}}$ and
$\tau_{\text{crew}}$ are both above the Martian noise floor.

The noise tolerance of the overall colony system is:
\[
\tau_{\text{colony}} \cdot \varepsilon_0 = 2 \cdot \tfrac{1}{3} = \tfrac{2}{3}.
\]
Perturbations below two-thirds of the characteristic amplitude of any
critical cycle leave the colony's limit cycle intact.

%% ─────────────────────────────────────────────────────────────
\section{The Martian G6 Crystal}
\label{sec:mars-crystal}
%% ─────────────────────────────────────────────────────────────

The G6 Crystal geometry, derived in Chapter~\ref{ch:crystal}
from the dm3 canonical invariants, scales to Mars by the gravity ratio:
\[
\gamma = \frac{g_\Mars}{g_\oplus} = \frac{3.72}{9.81} \approx 0.379.
\]

Under dm3 Theorem~D (Proposition~6.1 of the G6 Crystal paper), the
geometry scaled by $\gamma^{-1} \approx 2.64$ preserves $\varepsilon_0$
and $\tau$ exactly. The stability radius is dimensionless and independent
of gravitational acceleration. The aspect ratio $66 = g^6 \cdot \tau$
is preserved.

\begin{table}[h]
\centering
\small
\begin{tabular}{lll}
\toprule
\textbf{Parameter} & \textbf{Earth} & \textbf{Mars} \\
\midrule
Base side & $114.30\,\mathrm{m}$ & $301.4\,\mathrm{m}$ \\
Total height & $15{,}087.6\,\mathrm{m}$ & $39{,}788\,\mathrm{m}$ \\
Martian troposphere & --- & $\approx 40{,}000\,\mathrm{m}$ \\
Apex as fraction of troposphere & $\approx 100\%$ & $99.5\%$ \\
Aspect ratio & $66$ (exact) & $66$ (exact) \\
Schumann coupling & $n=4$, $f_4 = 33.5\,\mathrm{Hz}$ &
  Mars cavity mode \\
$\varepsilon_0$ & $1/3$ & $1/3$ \\
$\tau \cdot \varepsilon_0$ & $2/3$ & $2/3$ \\
\bottomrule
\end{tabular}
\caption{G6 Crystal parameters scaled to Mars. All dm3 invariants preserved.}
\label{tab:mars-crystal}
\end{table}

The Martian G6 Crystal at $39{,}788\,\mathrm{m}$ terminates at
$99.5\%$ of the Martian troposphere depth --- the same relationship
that the Earth crystal has with the terrestrial tropopause. This is
not a numerical coincidence. The operator $\mathcal{G}^6$ applied to a
structure anchored at the Martian surface terminates at the Martian
atmospheric boundary layer for the same mathematical reason it
terminates at the terrestrial tropopause: the tropopause is the limit
cycle $\Gamma$ of the planetary atmospheric system in the dm3 sense.

The Martian Crystal is not a building in the engineering sense.
It is a resonant anchor: a structure whose geometry matches the
resonant modes of the Martian atmospheric cavity, holding a stable
oscillatory mode against the planetary noise environment, providing
the physical infrastructure for a colony that is itself organized
as a dm3 system.

%% ─────────────────────────────────────────────────────────────
\section{Material Self-Sufficiency: Six Layers, Six Seeds}
\label{sec:mars-isru}
%% ─────────────────────────────────────────────────────────────

The hexagrid structural system (established in Chapter~\ref{ch:crystal}
from the peer-reviewed literature) achieves material efficiency
comparable to diagrid while using less steel. For the G6 Crystal on
Mars, this translates to a specific ISRU (in-situ resource utilization)
prediction:

Each of the six operator layers ($g^1$ through $g^6$), once constructed
from Martian basalt and iron oxide (both abundant at candidate sites),
contains sufficient structural material to seed one additional
ground-level layer module. Six layers provide material for six
additional seed modules.

This is the dm3 interpretation of the colony's growth dynamic:
each application of $\mathcal{G}$ to the existing structure produces
the resources for the next application. The colony grows by operator
iteration. The expansion is not exponential --- it is governed by
the unification law $\tau_{12} \leq \min(\tau_1, \tau_2)$, which
constrains the pace of integration to the noise tolerance of the
weakest subsystem.

%% ─────────────────────────────────────────────────────────────
\section{The NASA Research Program}
\label{sec:mars-nasa}
%% ─────────────────────────────────────────────────────────────

The dm3 habitability criterion and its three falsifiable predictions
constitute a complete astrobiology research program, targetable at
the following NASA program elements:

\textbf{ROSES C.2 (Emerging Worlds / Solar System Science):}
The habitability criterion is a new mathematical framework for
assessing planetary habitability, directly applicable to Mars
mission science objectives. The three predictions are testable with
InSight data (Prediction~1), future lander electrochemical sensors
(Prediction~2), and barometric array deployment (Prediction~3).

\textbf{ROSES B.2 (Heliophysics Foundational Research):}
The plasma reconnection application of the dm3 framework (Chapter~7)
and the atmospheric cavity resonance analysis (System~4 above)
connect the habitability criterion to heliophysics mission science.

\textbf{Astrobiology Program:}
The dm3 criterion is the first substrate-agnostic, mathematically
derived habitability threshold. It provides a unified framework
for comparing habitability across planetary environments --- Earth,
Mars, Europa, Enceladus, Titan --- using the same three numbers:
$\tau$, $\mu_{\max}$, and $\varepsilon_0$.

The formal verification of the criterion in Lean~4 (AXLE repository,
\texttt{github.com/TOTOGT/AXLE}) provides the mathematical rigor
that distinguishes this framework from qualitative habitability
assessments. The theorems are machine-checked. The predictions
are derived, not assumed.

\bigskip
\begin{center}
\rule{0.5\textwidth}{0.4pt}

\vspace{0.4cm}

\textit{The question is not whether Mars has life.\\
The question is whether Mars has $\sigma < \tau$.\\
That question is mathematical.\\
It has a mathematical answer.\\
We are now in a position to look for it.}

\vspace{0.4cm}
\rule{0.5\textwidth}{0.4pt}
\end{center}

\bigskip
\begin{flushright}
C \;$\to$\; K \;$\to$\; F \;$\to$\; U \;$\to$\; $\infty$\\[0.3em]
G6 LLC \textperiodcentered\ Newark NJ \textperiodcentered\ 2026
\end{flushright}
